% !PRINTED 361
\weberpage{361}
%<<CONT
so folgt daraus die Identität
\[
x_1 y_1 (x_2 y_2 - x_3 y_3) = (u_1 - u_2)(u_1 + u_2),
\]
und daraus wie oben
\[
\begin{aligned}
u_1 - u_2 &= h x_1 y_1\\
u_1 + u_2 &= \frac{x_2 y_2 - x_3 y_3}{h},
\end{aligned}
\]
worin $h$ eine Constante bedeutet, also
\[
2u_1 = h x_1 y_1 + \frac{x_2 y_2 - x_3 y_3}{h},
\]
und danach wird
\weq{7}{4f = 4 x_1 y_1 x_2 y_2 - \left(h x_1 y_1 + \frac{x_2 y_2 - x_3 y_3}{h}\right)^2.}
Dieser Gleichungsform können wir eine elegantere Gestalt geben, wenn wir
\[
h x_1,\quad \frac{x_2}{h},\quad \frac{x_3}{h}
\]
durch
\[
x_1,\quad x_2,\quad x_3
\]
ersetzen. Dann bedeuten die neuen $x_1, x_2, x_3$, gleich Null gesetzt, dieselben Linien wie die ursprünglichen, da sich ja beide nur durch einen constanten Factor unterscheiden, und es ergiebt sich aus (7):
\weq{8}{\begin{aligned}
-4f &= (x_1 y_1 + x_2 y_2 - x_3 y_3)^2 - 4 x_1 y_1 x_2 y_2\\
&= x_1^2 y_1^2 + x_2^2 y_2^2 + x_3^2 y_3^2\\
&\quad - 2 x_2 y_2 x_3 y_3 - 2 x_1 y_1 x_2 y_2 - 2 x_1 y_1 x_3 y_3,
\end{aligned}}
und die Gleichung $f = 0$ kann auch in der eleganten irrationalen Form
\[
x_1 y_1 + x_2 y_2 - x_3 y_3 = -2\sqrt{x_1 y_1 x_2 y_2}
\]
oder
\weq{9}{\sqrt{x_1 y_1} + \sqrt{x_2 y_2} + \sqrt{x_3 y_3} = 0}
dargestellt werden. Aus (9) können wir wieder rückwärts die Form (8) herleiten. Weil aber (9) ganz symmetrisch ist, so können wir die drei Paare vertauschen und erhalten z. B. auch
\[
4 x_2 y_2 x_3 y_3 = (-x_1 y_1 + x_2 y_2 + x_3 y_3)^2,
\]
woraus zu ersehen ist, dass auch die Berührungspunkte von $x_2 y_2 x_3 y_3$ auf einem Kegelschnitte liegen, und dass in dem Complex, den man aus $x_2 y_2$ erhält, nicht nur das Paar $x_1 y_1$, son-
%CONT>>
